\documentclass[10pt,letterpaper]{article}
\usepackage{technical-testing}
\setDocTitle{Threat Modeling Procedure}
\setDocSubtitle{Repeatable defensive engineering method}
\setDocVersion{1.0}
\setDocStatus{Independent engineering guidance}
\setDocOwner{latex-docs / security}
\setDocDate{2026-11-02}
\setDocSummary{Independent bounded procedure with verification and evidence criteria.}
\begin{document}
\makedoctocpage
\begin{notebox}[Source boundary]
Textbook concepts are anchored to Matt Bishop, \emph{Computer Security: Art and Science}, Second Edition (2019). All procedures, implementation examples, scoring rules and modern framework mappings below are independently authored supplemental engineering material, not excerpts, official compliance guidance or evidence of deployed controls.
\end{notebox}
\section{Required inputs and authority}
Collect current architecture, deployment inventory, data classification, identity flows, known exceptions, asset owners and scoped test authorization. Use synthetic examples when training. Do not record operational secrets in the public repository.
\section{Procedure}
\begin{enumerate}
\item \textbf{Define system and decision.} Record revision, in/out boundaries, users, trust dependencies, important outcomes and review owner.
\item \textbf{Draw a data flow diagram.} Number services, external actors, stores, cross-boundary flows and privilege changes; label authentication/authorization enforcement.
\item \textbf{Identify threats.} For each boundary consider identity spoofing, tampering, repudiation, disclosure, service exhaustion and privilege escalation (STRIDE is modern supplemental vocabulary).
\item \textbf{Write threat statements.} Actor, capability, entry point, exploitable condition, target asset, impact, observable evidence.
\item \textbf{Model abuse cases.} Include legitimate and malicious sequences plus state-dependent authorization conditions and non-obvious information flow.
\item \textbf{Prioritize and treat.} Owner decides impact and risk using documented method; map prevention, detection, containment and recovery mechanisms to assumptions they break.
\item \textbf{Verify.} Each material mitigation has a negative test and evidence artifact; unverified mitigations remain proposals.
\item \textbf{Revisit.} Trigger review on authentication, secrets, permissions, data model, deployment, network or third-party changes.
\end{enumerate}
\section{Minimum evidence and gate}
Produce Trust Boundary Diagram, Threat Register, Traceability record, Negative Tests, Residual-Risk Approval. Failure to identify ownership or review a critical externally exposed flow is a blocker under this proposed local policy.
\section{Source integration}
See existing \texttt{src/security/security-engineering/01-security-foundations-threat-modeling.tex} and \texttt{src/security/application-security/processes/}; this procedure adds book-specific formal-property checks instead of replacing those references.
\end{document}
